% Template: Sơ đồ flowchart cơ bản
% Dùng cho: giải thuật, quy trình, lưu đồ
% --------------------------------------------------------
\documentclass[border=10pt]{standalone}
\usepackage[utf8]{inputenc}
\usepackage[T5]{fontenc}
\usepackage[vietnamese]{babel}
\usepackage{tikz}
\usepackage{helvet}
\renewcommand{\familydefault}{\sfdefault}
% ── Hệ màu tập trung (chỉnh tại đây để đổi theme) ──────────
\usepackage{xcolor}
\definecolor{mauchinh} {HTML}{1565C0}
\definecolor{mauphu}   {HTML}{43A047}
\definecolor{maunoibat}{HTML}{E53935}
\definecolor{mauvien}  {HTML}{424242}
{mauchinh} {HTML}{1565C0}
{mauphu}   {HTML}{43A047}
{maunoibat}{HTML}{E53935}
{mauvien}  {HTML}{424242}
\usetikzlibrary{shapes.geometric, arrows.meta, calc}

\begin{document}
\begin{tikzpicture}[
    >=Stealth,
    line join=round,
    line cap=round,
    every node/.style={font=\sffamily\small},
    node distance=1.5cm,
    % ── Định nghĩa style các khối ──────────────────────
    start/.style={
        draw, rounded corners=10pt,
        fill=mauchinh!15, thick,
        minimum width=3cm, minimum height=0.8cm,
        text centered
    },
    process/.style={
        draw, rectangle,
        fill=white, thick,
        minimum width=3cm, minimum height=0.8cm,
        text centered
    },
    decision/.style={
        draw, diamond,
        fill=mauphu!20, thick,
        minimum width=2.5cm, minimum height=1cm,
        text centered, aspect=2
    },
    io/.style={
        draw, trapezium,
        trapezium left angle=70, trapezium right angle=110,
        fill=mauphu!10, thick,
        minimum width=3cm, minimum height=0.8cm,
        text centered
    },
    arrow/.style={thick, ->}
]
    % ── Các khối ─────────────────────────────────────────
    \path (0, 0)    node[start]    (start)   {Bắt đầu};
    \path (0,-1.8)  node[io]       (input)   {Nhập dữ liệu};
    \path (0,-3.6)  node[process]  (proc)    {Xử lý};
    \path (0,-5.4)  node[decision] (cond)    {Điều kiện?};
    \path (3,-5.4)  node[process]  (yes)     {Xử lý (Đúng)};
    \path (0,-7.2)  node[io]       (output)  {Xuất kết quả};
    \path (0,-9.0)  node[start]    (stop)    {Kết thúc};

    % ── Mũi tên nối ──────────────────────────────────────
    \draw[arrow] (start)  -- (input);    % Bắt đầu → Nhập
    \draw[arrow] (input)  -- (proc);     % Nhập → Xử lý
    \draw[arrow] (proc)   -- (cond);     % Xử lý → Điều kiện
    \draw[arrow] (cond)   -- node[above] {Đúng} (yes);    % → Nhánh đúng
    \draw[arrow] (cond)   -- node[right] {Sai}  (output); % → Xuất kết quả
    \draw[arrow] (yes)    |- (output);   % Nhánh đúng → Xuất
    \draw[arrow] (output) -- (stop);     % Xuất → Kết thúc
\end{tikzpicture}
\end{document}
